% Generated draft: requires mathematical and source review.
\documentclass[11pt]{article}
\usepackage[margin=25mm]{geometry}
\usepackage{fontspec}
\setmainfont{Latin Modern Roman}
\usepackage{amsmath,amssymb,mathtools,unicode-math}
\setmathfont{Latin Modern Math}
\usepackage{enumitem,needspace,xurl,hyperref}
\setlength{\parindent}{0pt}
\setlength{\parskip}{5pt}
\setlength{\emergencystretch}{3em}
\newcommand{\N}{\mathbb{N}}
\newcommand{\Z}{\mathbb{Z}}
\newcommand{\Q}{\mathbb{Q}}
\newcommand{\R}{\mathbb{R}}
\newcommand{\C}{\mathbb{C}}
\newcommand{\F}{\mathbb{F}}
\newcommand{\A}{\mathbb{A}}
\newcommand{\PP}{\mathbb{P}}
\newcommand{\E}{\mathbb{E}}
\newcommand{\eps}{\varepsilon}
\newcommand{\dd}{\,\mathrm{d}}
\begin{document}
\section*{UnsolvedMath ID: 30003301}
\textbf{Source problem code:} catalogue entry
\subsection{Definitions and mathematical statement}
Under the conventional interpretation of the notation, let \(S_g\) be a closed, connected, oriented surface of genus \(g\), and let \(p\in S_g\) be a distinguished point. Write \(\operatorname{Mod}_g\) for the group of orientation-preserving mapping classes of \(S_g\), and \(\operatorname{Mod}_{g,1}\) for the group of orientation-preserving mapping classes fixing \(p\), with isotopies also fixing \(p\). Thus the latter is interpreted as a surface with one marked point, equivalently one puncture, rather than one boundary component. Forgetting the marked point gives a homomorphism
\[
F:\operatorname{Mod}_{g,1}\longrightarrow\operatorname{Mod}_g.
\]
For a simple closed curve \(a\), let \(t_a\) denote its Dehn twist, using a fixed handedness convention throughout.

\textbf{Question.} Given a finite ordered collection of simple closed curves \(a_1,\ldots,a_n\) on \(S_g\) satisfying
\[
t_{a_1}\cdots t_{a_n}=1\qquad\text{in }\operatorname{Mod}_g,
\]
do there exist simple closed curves \(\widetilde a_1,\ldots,\widetilde a_n\subset S_g\setminus\{p\}\), each isotopic to the corresponding \(a_i\) after forgetting \(p\), such that
\[
t_{\widetilde a_1}\cdots t_{\widetilde a_n}=1\qquad\text{in }\operatorname{Mod}_{g,1}?
\]
In particular, each chosen twist must satisfy \(F(t_{\widetilde a_i})=t_{a_i}\), and the ordered product itself must be the identity before forgetting the puncture.

\textbf{Unresolved qualifications.} The source supplies no range for \(g\), no explicit essentiality or nonseparating condition on the curves, and no handedness convention. The displayed source relation contains twists without inverse exponents. The marked-point interpretation and the interpretation of ``moved'' as individual isotopy after filling the puncture make the question precise, but require confirmation against the original report. The source does not explicitly resolve whether a more restrictive notion of moving the entire curve configuration is intended.
\subsection{Short English statement}
Can every supplied ordered relation among Dehn twists on a closed surface be realized by corresponding curves avoiding one puncture, with the same ordered product equal to the identity in the punctured mapping class group? The genus range, curve restrictions, and precise meaning of moving the curves require source confirmation. The supplied record labels the problem open; this draft does not independently verify that status.
\subsection{Sources}
\begin{itemize}
\item Supplied local UnsolvedMath registry record, permanent ID \(30003301\), titled ``Lifting Dehn-Twist Relations to Punctured Surfaces,'' source problem code \texttt{catalogue entry}, in the Oberwolfach Reports Open Problems collection. No individual proposer is supplied; the proposed year is recorded as 2016.
\item Supplied report citation: \emph{Surface Bundles} (2017), Oberwolfach Reports, pp.\ 3149--3195, DOI \texttt{10.4171/owr/2016/56}. \url{https://doi.org/10.4171/owr/2016/56}. The literature-source label is ``OWR 56 (2017).'' The supplied background describes these reports as documenting mathematical workshops and their open-problem discussions.
\item Supplied literature reference: Gervais, ``a presentation of mapping class groups.'' \url{https://arxiv.org/abs/math/9801025}.
\item Literature assessment checked at \texttt{2026-08-21}: the supplied curation labels the problem open and reports that no checked primary source gives a general theorem lifting every Dehn-twist relation from \(\operatorname{Mod}_g\) to curves avoiding the puncture in \(\operatorname{Mod}_{g,1}\). This is supplied dated curation derived from the Oberwolfach Reports record, not an independently verified literature assessment. The background explicitly requests independent verification of time-sensitive status and cited claims.
\end{itemize}
\end{document}
